\section{Attack C: Shai-Hulud 1.0}
\label{sec:attackC}

On September 15, 2025, 187 \ac{NPM} packages, 
including @ctrl/tinycolor with over two million weekly downloads, 
were infected with malicious code designed to steal 
sensitive information such as tokens and 
authentication keys \cite{oxsecurity2025npm20hack}.

The method of the first infection and the patient zero remained unknown,
but it could be a case of phishing \cite{kaspersky2025tinycolorshaikuludsupplychainattack}.
The number of infected packages is so high 
because the malware automatically publishes 
a malicious version of every package of the compromised maintainer (of which it has publishing permissions).

The malicious code was inserted into the packages' tarballs
in the new file bundle.js.

The file weighs approximately 3.7 MB and contains only a single line of code
3,734,355 characters long.

When the victim finishes installing the package with the infected version,
the malicious script is executed and downloads a platform-appropriate version of TruffleHog,
a legitimate tool for searching secrets such as cryptographic keys and \ac{API} tokens
in local file systems and Git repositories.

After starting and completing the search with TruffleHog, 
the malicious code publishes the stolen data 
by creating a public repository named Shai-Hulud 
(like the worm in the Dune novels) 
on behalf of the victim using their GitHub token.

In this repository, there is a file data.json containing all
the collected secrets and system information encoded.